\begin{titlepage}
\thispagestyle{empty}
\begin{tikzpicture}[remember picture, overlay]
  \fill[deepblue] (current page.north west) rectangle (current page.south east);
  \fill[softcopper!18] ([xshift=-0.7cm,yshift=-2.1cm]current page.north east) circle (6.5cm);
  \fill[white!8] ([xshift=1.3cm,yshift=2.1cm]current page.south west) circle (5.0cm);

  \draw[softcopper!80, line width=1.4pt]
    ([xshift=1.1cm,yshift=-1.1cm]current page.north west)
    rectangle
    ([xshift=-1.1cm,yshift=1.1cm]current page.south east);

  \node[
    anchor=north west,
    fill=softcopper,
    text=deepblue,
    font=\bfseries\large,
    rounded corners=4pt,
    inner xsep=10pt,
    inner ysep=6pt
  ] at ([xshift=1.6cm,yshift=-1.5cm]current page.north west) {QUYỂN XV};

  \node[
    anchor=north,
    text=white,
    font=\large\itshape
  ] at ([yshift=-1.7cm]current page.north) {Truyện Tích Cảm Hứng Song Ngữ};

  \node[
    anchor=west,
    text=white,
    font=\fontsize{22}{26}\selectfont\bfseries
  ] at ([xshift=1.8cm,yshift=4.2cm]current page.center) {NHỮNG ĐIỀU};

  \node[
    anchor=west,
    text=softcopper,
    font=\fontsize{22}{26}\selectfont\bfseries
  ] at ([xshift=1.8cm,yshift=2.3cm]current page.center) {THẦY MUỐN NÓI};

  \node[
    anchor=west,
    text=white,
    font=\fontsize{22}{26}\selectfont\bfseries
  ] at ([xshift=1.8cm,yshift=0.4cm]current page.center) {VỚI HỌC SINH};

  \draw[softcopper, line width=2pt]
    ([xshift=1.8cm,yshift=-1.0cm]current page.center) --
    ([xshift=10.1cm,yshift=-1.0cm]current page.center);

  \node[
    anchor=west,
    text=white,
    font=\large
  ] at ([xshift=1.8cm,yshift=-2.1cm]current page.center) {Một bộ sách thầy nói chuyện thẳng, gần và chậm với học sinh về học tập, tính cách và đường đời};

  \node[
    anchor=north west,
    fill=white,
    text=black,
    rounded corners=8pt,
    text width=11.7cm,
    inner sep=14pt,
    align=left,
    font=\normalsize
  ] at ([xshift=1.9cm,yshift=-14.9cm]current page.north west) {%
    \textbf{Tinh thần quyển này:} không giảng đạo lý theo kiểu xa cách. Mỗi bài viết là một lời nói gần gũi, thẳng nhưng dịu, để học sinh thấy mình được nói chuyện chứ không chỉ bị dạy bảo.\\[6pt]
    \textbf{Cách dùng:} đọc từng bài ngắn mỗi ngày, học thêm vài câu tiếng Anh đơn giản, và dùng để gợi mở giờ sinh hoạt, chủ nhiệm hoặc tự học.\\[6pt]
    \textbf{Điểm tựa:} gần lớp học, gần tâm lý tuổi học trò, dễ đọc, dễ nhớ, và giữ nhịp đồng bộ với cả bộ sách.
  };

  \node[
    anchor=south,
    text=softcopper,
    font=\large\itshape,
    align=center
  ] at ([yshift=2.1cm]current page.south) {%
    ``Có những điều học sinh cần nghe không phải vì các em không biết,\\
    mà vì các em cần một người nói lại bằng giọng đủ gần.''%
  };

  \node[
    anchor=south,
    text=white!85,
    font=\normalsize,
    align=center
  ] at ([yshift=1.1cm]current page.south) {Dành cho học sinh, giáo viên và những ai muốn nói với người trẻ bằng một giọng bình tĩnh hơn};
\end{tikzpicture}
\end{titlepage}
\clearpage
